% Lösungen Einheit 3 von 3 – Kugel (zusammenbau v0.1)
\einheitenkopf{Einheit 3 von 3 · Kugel}

%% TODO Lösung der Erklärzeile 21g) fehlt in der Bank
\erg{21}{a) Radius \quad b) $V = \frac{4}{3} \cdot \pi \cdot 3^3 \approx 113{,}1$ cm³ \quad c) $V = \frac{4}{3} \cdot \pi \cdot 4^3 \approx 268{,}1$ cm³ \quad d) $V = \frac{4}{3} \cdot \pi \cdot 2^3 \approx 33{,}5$ dm³ \quad e) $V = \frac{4}{3} \cdot \pi \cdot 7^3 \approx 1\,436{,}8$ cm³ \quad f) $V = \frac{4}{3} \cdot \pi \cdot 9^3 \approx 3\,053{,}6$ m³ \quad g) \ldots \quad h) $r = 8$ cm; $V = \frac{4}{3} \cdot \pi \cdot 8^3 \approx 2\,144{,}7$ cm³ \quad i) $V = \frac{4}{3} \cdot \pi \cdot 2{,}5^3 \approx 65{,}4$ cm³ \quad j) $V = \frac{4}{3} \cdot \pi \cdot 15^3 \approx 14\,137{,}2$ cm³ $\approx 14{,}1$ l \quad k) $O = 4 \cdot \pi \cdot 11^2 \approx 1\,520{,}5$ cm² \quad l) Kreis mit dem Radius $2$ cm; flache Ellipse durch den Mittelpunkt als Äquator (hintere Hälfte gestrichelt); Radius vom Mittelpunkt zum Rand gestrichelt, mit $2$ cm beschriftet \quad m) $V = \frac{1}{2} \cdot \frac{4}{3} \cdot \pi \cdot 4^3 \approx 134{,}0$ cm³; $O = 2 \cdot \pi \cdot 4^2 + \pi \cdot 4^2 \approx 150{,}8$ cm² \quad n) $V = \frac{4}{3} \cdot \pi \cdot 2^3 \approx 33{,}51$ cm³; $m \approx 33{,}51 \cdot 7{,}9 \approx 265$ g \quad o) $r = \sqrt{200 : (4 \cdot \pi)} \approx 3{,}99$ cm \quad p) Kante $= 28$ cm; Würfel $= 21\,952$ cm³; Kugel $\approx 11\,494{,}0$ cm³; Luft $\approx 10\,458{,}0$ cm³ \quad q) Zylinder $\pi \cdot 3^2 \cdot 9 \approx 254{,}5$ m³; Halbkugel $\frac{2}{3} \cdot \pi \cdot 3^3 \approx 56{,}5$ m³; zusammen $\approx 311{,}0$ m³ \quad r) A: $V \approx 33{,}5$ cm³, B: $V \approx 268{,}1$ cm³; $268{,}1 : 33{,}5 \approx 8$ – das Achtfache, weil $2^3 = 8$ \quad s) $r = 10{,}8$ cm; $V = \frac{4}{3} \cdot \pi \cdot 10{,}8^3 \approx 5\,276{,}7$ cm³}

21l)

\kugel{2}{$2$ cm}

\erg{22}{a) Nora hat den Durchmesser eingesetzt; der Radius ist $9$ cm. Richtig: $V = \frac{4}{3} \cdot \pi \cdot 9^3 \approx 3\,053{,}6$ cm³ \quad b) Zur halben gewölbten Fläche kommt die ebene Schnittfläche, ein ganzer Kreis, dazu. \quad c) Zylinder mit dem Radius $2{,}5$ m und der Höhe $8$ m, darauf eine Halbkugel mit dem Radius $2{,}5$ m; Gesamthöhe $= 8 + 2{,}5 = 10{,}5$ m \quad d) gewölbte Fläche: $2 \cdot \pi \cdot 2{,}8^2 \approx 49{,}26$ m² (ohne Schnittkreis); zweimal $\approx 98{,}52$ m²; $98{,}52 : 15 \approx 6{,}57$, aufgerundet $7$ Eimer}
